% =====================================================================
%  9. Conclusão
% =====================================================================
\section{Conclusão}
\label{sec:conclusao}

\indent\indent Este trabalho descreveu dez métodos de avaliação de frases isoladas em português sob um
formalismo único e os submeteu ao mesmo conjunto controlado de variantes, distribuídos por
quatro dimensões: ortografia, gramática, estrutura e semântica.

A caracterização não produziu uma ordenação. Nenhum dos dez domina os demais, e as falhas
observadas não se corrigem umas às outras porque nascem em níveis distintos: umas na definição
do próprio método, outras na cobertura finita do recurso que ele consulta, outras ainda no
insumo que um modelo anterior lhe entrega.

Sendo assim decorre que combinar os métodos não é conveniência de implementação, mas a resposta que o
levantamento indica. Como as coberturas são complementares e as falhas não coincidem, o que
um método deixa passar é frequentemente o que outro alcança --- e a cadeia proposta reúne
quatro etapas, uma por dimensão, decidindo pela existência de achados localizados em lugar de
escore agregado. Cada item do laudo conserva a etapa que o produziu e as palavras envolvidas, e
um único limiar figura em toda a cadeia.

Os limites são de duas ordens, e ambas recomendam cautela na leitura. O corpus é controlado e
pequeno, derivado de uma única frase de referência, e cada método foi aplicado às variantes que
exercitam a propriedade que ele mede, de modo que a complementaridade permanece argumento de
construção mais do que medida. E os instrumentos são únicos --- um analisador sintático, um
corpus de contagens, um modelo mascarado ---, nenhum treinado ou ajustado aqui, o que significa
que parte do que se observou pertence a eles, e não às definições.

Desses limites decorrem os trabalhos futuros. O primeiro é completar a aplicação cruzada,
submetendo cada método a todas as variantes, para que a complementaridade deixe de ser
argumento e passe a ser medida. O segundo é ampliar o corpus, tanto em número de frases de
referência quanto em tipos de desvio, e submetê-lo a julgamentos humanos de aceitabilidade,
de que este trabalho não dispôs. O terceiro é variar os instrumentos, repetindo as medições sob
outro analisador, outro corpus de contagens e outro modelo mascarado, o que permitiria separar
o que pertence às definições do que pertence às ferramentas. O quarto é incorporar as regras
escritas à mão à cadeia, uma vez que alcançam a escolha lexical, desvio que as quatro etapas
atuais atravessam sem detectar. Resta, por fim, o erro factual, que é português correto e falso
apenas em relação à base que originou o texto: sua verificação exige confrontar o enunciado com
os campos que o geraram, e devolve o problema ao sistema gerador.
